\chapter{Resumen}
La industria marítima mundial está dando cada vez más prioridad a la sostenibilidad, promoviendo la adopción de diversos sistemas de propulsión híbridos y totalmente eléctricos para diferentes tipos de buques. En este artículo se presenta una metodología de diseño y dimensionamiento de un sistema de almacenamiento de energía para un remolcador híbrido, personalizado para satisfacer las exigencias operativas y minimizar el impacto medioambiental. El marco propuesto integra el análisis del perfil de carga, el dimensionamiento del almacenamiento de energía, la selección de tecnologías de baterías y la optimización del sistema de propulsión para lograr un equilibrio óptimo entre el rendimiento y la reducción de emisiones. Los resultados demuestran que un sistema de almacenamiento de energía bien diseñado para un remolcador híbrido puede lograr reducciones en el consumo de combustible y las emisiones de gases de efecto invernadero sin comprometer la disponibilidad y la robustez necesarias para las operaciones de remolque y tránsito.


